% !TeX root = ../sections/02_exercises.tex

\begin{exercise}
		R-2. 環
		\[
		\mathbb{Z}[\sqrt{2}+\sqrt{3}]
		=
		\{\, f(\sqrt{2}+\sqrt{3}) \mid f(X)\in\mathbb{Z}[X] \,\}
		\]
		は，$\mathbb{Z}$-加群とみたとき $\mathbb{Z}$-自由加群であることを示せ。
		また，その階数も求めよ。
	\end{exercise}
	
	\begin{background}
		$\alpha=\sqrt{2}+\sqrt{3}$ とおく。
		このとき $\mathbb{Z}[\alpha]$ は，
		$\alpha$ の整数係数多項式で表される元全体である。
		
		$\alpha$ が有理数体 $\mathbb{Q}$ 上代数的で，その最小多項式の次数が $d$ であれば，
		\[
		1,\alpha,\alpha^2,\dots,\alpha^{d-1}
		\]
		が $\mathbb{Q}(\alpha)$ の $\mathbb{Q}$-基底になる。
		
		今回の $\alpha=\sqrt{2}+\sqrt{3}$ については，
		\[
		\alpha^2=5+2\sqrt{6}
		\]
		であるから，
		\[
		(\alpha^2-5)^2=24
		\]
		となる。したがって
		\[
		\alpha^4-10\alpha^2+1=0
		\]
		が成り立つ。
		
		つまり，$\alpha$ は
		\[
		X^4-10X^2+1
		\]
		の根である。
	\end{background}
	
	\begin{strategy}
		まず
		\[
		\alpha=\sqrt{2}+\sqrt{3}
		\]
		とおく。
		関係式
		\[
		\alpha^4-10\alpha^2+1=0
		\]
		より，$\alpha^4$ は $1,\alpha,\alpha^2,\alpha^3$ の
		$\mathbb{Z}$-係数線形結合で表せる。
		したがって，$\alpha^4$ 以上の高次の冪はすべて
		\[
		1,\alpha,\alpha^2,\alpha^3
		\]
		の $\mathbb{Z}$-係数線形結合で表せる。
		
		よって，$\mathbb{Z}[\alpha]$ は
		\[
		1,\alpha,\alpha^2,\alpha^3
		\]
		で生成される。
		
		次に，これらが $\mathbb{Z}$ 上一次独立であることを示す。
		そのためには，$\alpha$ の最小多項式が
		\[
		X^4-10X^2+1
		\]
		であり，次数が $4$ であることを確認すればよい。
		
		以上により，
		\[
		\{1,\alpha,\alpha^2,\alpha^3\}
		\]
		が $\mathbb{Z}[\alpha]$ の $\mathbb{Z}$-基底となる。
		したがって，
		\[
		\mathbb{Z}[\sqrt{2}+\sqrt{3}]
		\]
		は階数 $4$ の $\mathbb{Z}$-自由加群である。
	\end{strategy}